\section{Closing the remaining higher symplectic source in characteristic three}
\label{sec:symplectic-three-semisimple}

\paragraph{Status and notation.}
This separate complete new derivation has passed whole-proof reviews
by /root and /root/imprimitive, with separate exact supplements. It uses D-HOS-17--18, leaving their frozen proof
unchanged. The finite checks described below supplement, rather than
replace, the universal arguments. Write $S_n=\operatorname{Sp}_{2n}(3)$
and $\nu(T)$ for the codimension of its largest complex eigenspace.

\begin{lemma}[D-HOS-20: odd-prime fixed forms]
\label{D-HOS-20}
Let $A\in\operatorname{GL}_{N}(3)$ have prime order $\ell\ge5$
and have no eigenvalue $1$. For every nonzero general-linear
congruence orbit $\mathcal O$ of symmetric bilinear forms,
\[
             |\{B\in\mathcal O:A^tBA=B\}|\le|\mathcal O|/9.
                                                               \tag{HOS3P.1}
\]
\end{lemma}

\begin{proof}
The matrix $A$ is semisimple, and it has neither $1$ nor $-1$ as
an eigenvalue. A fixed form has an $A$-stable radical. On its
nondegenerate quotient of dimension $r>0$, the induced matrix
$A_0$ still has order $\ell$ and no such eigenvalue. Over the
algebraic closure, an invariant symmetric form pairs only
eigenspaces with eigenvalues $\lambda$ and $\lambda^{-1}$.
Nondegeneracy requires equal dimensions within each inverse pair.
Thus $r$ is even. If $r=2$, the matrix would belong to
$O_2^+(3)$ or $O_2^-(3)$, whose orders are four and eight,
respectively. This is impossible for $\ell\ge5$. These orders
follow directly from the diagonal and norm-one descriptions of
the two binary orthogonal groups, or from
\cite[Theorem~11.4 and its proof, pp.~139--140]{HOS:Taylor1992}.
Hence any nonzero fixed form has quotient dimension $r\ge4$.

Write $m_j$ for the common dimensions in the inverse eigenvalue
pairs of $A_0$. The vector space of all invariant symmetric forms
on the quotient has dimension
\[
                \sum_jm_j^2\le\left(\sum_jm_j\right)^2
                         =r^2/4.                            \tag{HOS3P.2}
\]
The dimension is the same over $\mathbf F_3$ and over its
algebraic closure because it is the kernel dimension of a linear
system over $\mathbf F_3$. It therefore contains at most
$3^{r^2/4}$ nondegenerate forms of either prescribed type.

For comparison, fix a nonsingular diagonal matrix $D$ of that
type. For every upper unitriangular matrix $L$, the matrix
$L^tDL$ has the same type. These matrices are distinct: if two
are equal, their quotient is an upper unitriangular isometry of
$D$; its first column is the first standard vector, so
orthogonality forces all other entries of its first row to vanish,
and induction gives the identity. Thus the type has at least
$3^{r(r-1)/2}$ nondegenerate forms. For each fixed radical,
the fixed proportion of quotient forms is at most
\[
             3^{r^2/4-r(r-1)/2}
                =3^{-r(r-2)/4}\le1/9.                       \tag{HOS3P.3}
\]
Every form in the original orbit is specified by its radical
and a nondegenerate quotient form of the specified type. The
number of quotient forms is independent of the chosen radical.
Non-$A$-stable radicals contribute no fixed forms; for each
stable radical the preceding estimate applies. Summing proves
(HOS3P.1). No enumeration of general-linear orbits is assumed.
\end{proof}

\begin{theorem}[D-HOS-21: all odd-prime semisimple elements]
\label{D-HOS-21}
Let $n\ge4$, let $V$ be a nontrivial ordinary irreducible complex
$S_n$-module of degree $a$, and let $h\in S_n$ have prime order
$\ell\ge5$. Then
\[
                            \nu(h|V)\ge4a/15.               \tag{HOS3P.4}
\]
This includes all characters inflated from represented central
quotients and requires no selfduality assumption.
\end{theorem}

\begin{proof}
If $h$ has a natural eigenvalue $1$, its fixed space is a
nondegenerate symplectic space, hence contains a hyperbolic pair.
The phase-tolerant Heisenberg argument already proved in
D-HOS-14 gives (HOS3P.4). Since $h$ has odd order, $-1$ cannot
be a natural eigenvalue. We therefore assume that it has neither.

We first treat any case in which there are opposite $h$-stable
totally isotropic subspaces $L,L'$ of common dimension $N\ge4$
whose sum $W_0$ is nondegenerate. The subgroup
$S_0=\operatorname{Sp}(W_0)$ acts identically on its perpendicular
space. Its two opposite Siegel radicals $U,U^-$ generate $S_0$
and are conjugate, by the explicit argument in D-HOS-18.
If $S_0=S_n$, its invariant subspace in $V$ is zero. Otherwise
Larsen--Shalev--Tiep's precise restriction estimate
\cite[Proposition~4.2.3 and complete proof, pp.~1909--1912]{HOS:LST2011}
gives
\[
 \frac{\dim V^{S_0}}a<3^{8/3-N}\le3^{-4/3}<1/4.
                                                               \tag{HOS3P.5}
\]
The subspace intersection inequality for $V^U,V^{U^-}$ therefore
gives $\dim V^U<5a/8$ (with the stronger $a/2$ bound if $S_0=S_n$).

The dual of $U$ is the symmetric-form space on $L$, with the
congruence action of its full Levi subgroup $\operatorname{GL}(L)$.
Consequently weight multiplicities are constant on each orbit.
The matrix $A=h|_L$ has order $\ell$ and no eigenvalue $1$.
D-HOS-20 bounds the fixed fraction in every nonzero orbit by
$1/9$. All other weight orbits of $h$ have length $\ell$.
An eigenvector on such an orbit is determined by its component
in one weight space, so each eigenvalue occupies at most a
$1/\ell$ fraction there. Thus
\[
 \nu(h|V)\ge(1-1/\ell)\frac89\,(a-\dim V^U)
             >\frac45\frac89\frac38a=\frac4{15}a.           \tag{HOS3P.6}
\]

It remains to show that the absence of suitable $L,L'$ forces
a small centralizer. We give the primary-space argument fully.
Put $f=\operatorname{ord}_{\ell}(3)$. Every nontrivial irreducible
factor $P$ of $X^\ell-1$ over $\mathbf F_3$ has degree $f$.
The primary space $W_P$ pairs under the symplectic form only
with $W_{P^*}$, where $P^*$ is the reciprocal polynomial. This
follows either by diagonalizing over the algebraic closure, or
by observing that the adjoint of $h$ is $h^{-1}$.

If $f$ is odd, no such $P$ is self-reciprocal: the condition
$\lambda^{-1}=\lambda^{3^j}$ for a primitive $\ell$-th root
would put the order-two element $-1$ in the odd-order subgroup
generated by $3$ modulo $\ell$. Choose one primary space from
each pair to form $L$ and the other to form $L'$. They are
opposite Lagrangians, of dimension $n\ge4$, so (HOS3P.6) applies.

If $f$ is even, write $f=2d$. The unique element of order two
in the group generated by $3$ modulo $\ell$ is
$3^d=-1$, so every primary polynomial is self-reciprocal.
Moreover $d\ge2$: degree two would require the prime $\ell\ge5$
to divide $3^2-1=8$. Put $E=\mathbf F_{3^{2d}}$ and
$E_0=\mathbf F_{3^d}$. On a primary space of multiplicity $m_P$,
$h$ acts as multiplication by a root $\lambda$ generating $E$,
with adjoint involution $\sigma(\lambda)=\lambda^{-1}$.

For clarity, the associated Hermitian structure follows directly
from trace duality. There is a unique $E$-valued pairing $K$ such
that
\[
 \operatorname{Tr}_{E/\mathbf F_3}(bK(x,y))
                         =\langle bx,y\rangle\quad(b\in E).
                                                               \tag{HOS3P.7}
\]
The adjoint rule for $E$ gives
$K(bx,y)=bK(x,y)$, $K(x,by)=\sigma(b)K(x,y)$ and
$K(y,x)=-\sigma K(x,y)$. It is nondegenerate. Choose a nonzero
$c\in E$ with $\sigma(c)=-c$; then $cK$ is Hermitian.
Conversely the original alternating form is recovered by the
trace in (HOS3P.7), so an $E$-isotropic Hermitian subspace is
also symplectically isotropic over $\mathbf F_3$.

The finite Hermitian Witt decomposition has index
$\lfloor m_P/2\rfloor$, as proved in
\cite[Lemma~10.2, Corollary~10.3 and the following decomposition,
pp.~116--117]{HOS:Taylor1992}.
If any $m_P\ge2$, one Hermitian hyperbolic plane supplies
opposite $\mathbf F_3$-isotropic spaces of dimension
$2d\ge4$, whose sum is symplectically nondegenerate. Again
(HOS3P.6) applies.

The only case left is therefore $m_P=1$ for every primary
polynomial. Distinct primary spaces cannot be permuted by an
operator commuting with $h$. On each space that operator is
$E$-linear, and hence multiplication by an element $b\in E^*$.
Preserving the alternating form, or equivalently (HOS3P.7),
is exactly $b\sigma(b)=1$. Thus the centralizer in the full
symplectic group has order
\[
               |C_{S_n}(h)|=(3^d+1)^k,\qquad n=dk.          \tag{HOS3P.8}
\]
The scalar center is already included in this product; no
projective-group order is substituted.

For all integers $d\ge2$, $3^d+1\le10^{d/2}$: equality holds
at $d=2$, and the ratio on increasing $d$ is less than
$3/\sqrt{10}<1$. Hence $|C_{S_n}(h)|\le10^{n/2}$.
Character column orthogonality gives
$|\chi(h)|\le|C_{S_n}(h)|^{1/2}$, the elementary statement
\cite[Lemma~4.3.1, p.~1912]{HOS:LST2011}.
Using the proven degree bound D-HOS-02 gives
\[
 \frac{|\chi(h)|}{a}
       \le\frac{10^{n/4}}{3^{n-1}}\le\frac{10}{27},
 \qquad
       \nu(h|V)\ge\frac{1-10/27}{2}a
                         =\frac{17}{54}a>\frac4{15}a.       \tag{HOS3P.9}
\]
The middle expression decreases with $n$, since $10^{1/4}<3$,
and its value at $n=4$ is $10/27$. The support inequality uses
$|\chi(h)|\ge a-2\nu(h|V)$ and remains valid without a
faithfulness assumption on the character. The cases are
exhaustive, proving (HOS3P.4).
\end{proof}

\begin{theorem}[D-HOS-22: the higher characteristic-three source is excluded]
\label{D-HOS-22}
For every $n\ge4$, every nontrivial ordinary irreducible complex
$S_n$-module of degree $a$, and every noncentral $g\in S_n$,
\[
 \nu(g|V)\ge\frac29a\quad(n\text{ even}),\qquad
 \nu(g|V)\ge\frac{19}{104}a\quad(n\text{ odd}).               \tag{HOS3P.10}
\]

A finite inducing image with an irreducibly acting normal central
product of such factors, all with nonselfdual factor modules,
cannot supply period rank nine for its actual six-dimensional
multiplicity variation. Different natural ranks, represented
central quotients, scalar extensions and all compatible factor
permutations are allowed. Consequently the transitive ordinary
symplectic-core case in the main imaginary-quadratic sixfold
problem has no source with $q=3$ and $n\ge4$, in either
determinant class.
\end{theorem}

\begin{proof}
The projective kernel of a nontrivial irreducible representation
of $S_n$ is its scalar center. Taking a noncentral power of $g$
of prime order modulo that center, and changing its central
sign if necessary, cannot increase projective support. For odd
prime order the representative has actual prime order; for
projective order two its square is $I$ or $-I$.

Actual order three has support at least $2a/9$ by D-HOS-13.
The case $h^2=I$ has that same bound by the already proved
Heisenberg argument of D-HOS-14. D-HOS-18 treats $h^2=-I$,
with bounds $19a/78>2a/9$ when $n$ is even and
$19a/104$ when $n$ is odd. D-HOS-21 treats every remaining
odd prime, with bound $4a/15>2a/9$. These exhaust all cases,
so (HOS3P.10) follows.

For a nonselfdual factor module the only compatible factor
normalizers are scalar times inner, by the exact class-inverting
diagonal-automorphism argument in D-HOS-15. Factorwise tensor
operators retain the support fraction of any nonscalar factor
by convolution. A nontrivial factor permutation has support at
least one quarter of the full dimension. A nonidentity scalar
has actual moving rank equal to that dimension.

If every factor has even $n_i$, every genuine inertia thus moves
at least $2e/9$ dimensions. If some factor has odd $n_i$, the
uniform bound $19e/104$ applies. The factor degree bound gives
$e\ge27$ in the first case and $e\ge81$ in the second.
As before the exact intermediate-cover Euler formula
$6=2e(h-1)+D$ excludes genus at least two; genus one has at most
six genuine marks and period-image dimension at most six.
Rank nine therefore requires genus zero, at least twelve
genuine inertia points, and $D=2e+6$. It would force
\[
                 e\le9\quad\hbox{or}\quad e\le156/5,
                                                               \tag{HOS3P.11}
\]
respectively, contradicting the preceding lower bounds.
No invariant bilinear form is required for this nonselfdual
normal-core conclusion, and the inducing character field is
not identified with the original imaginary quadratic field.

In a transitive core selfduality is preserved by automorphic
pullback, so either every factor module is selfdual or none is.
The former case is already excluded in the main imaginary-field
setting by D-HOS-12. The latter is excluded above. This proves
the final assertion; the small ranks $n<4$ remain within their
separately stated earlier results.
\end{proof}

\paragraph{Primary inputs and certificate scope.}
The new semisimple primary-space and centralizer reductions, the
unitriangular count and the source argument are proved here.
The published restriction estimate is used at its exact ordinary
simply connected-group scope, as checked in D-HOS-18. The small
binary orthogonal orders and finite Hermitian decomposition are
the stated elementary primary inputs from Taylor. The accompanying
exact finite supplement tests the fixed-form argument, both kinds
of primary behavior, a toral centralizer and all constants. It
does not infer the universal theorem from those finite examples.
